\section{Passage to a global estimate}\label{sec:global}

The packet theorem first gives a cubic estimate on widely separated
balls, with an arbitrarily small power loss.  We remove that loss by
covering a finite list of peaks with sparse families of balls.  The
initial finiteness of the peak list follows from an elementary
fourth-power estimate for the original sphere measure.

The choices have a fixed order: first the separation exponent \(C_0\),
then the desired distributional error \(\alpha\), and finally the
power loss \(\epsilon\).  The covering radii may depend on \(\alpha\);
the sparse estimate must therefore permit \(\epsilon\) to be chosen
after those radii have been bounded.

This passage from local estimates to a global bound follows the
sparse-family strategy of Tao~\cite{Tao1999Epsilon}; compare the
bounded-data formulation in
\cite[arXiv:1012.3760v3, Appendix, Lemmas~A1--A2]{BourgainGuth2011}.
We prove the required covering and cubic estimates below.

\subsection{A sparse estimate for graph patches}

Normalize first by $\|g\|_{L^\infty(S^2,\sigma)}=1$. Partition the
sphere into finitely many measurable pieces contained in the graph patches
of Section~\ref{sec:packets}. Subdivide these pieces, if necessary, so
that on each piece one of the terminal windows $\chi_1^{\nu_0}$ is
bounded below by a positive constant. Such a subdivision exists because
the windows form a square partition of unity with bounded overlap.
In the rotated coordinates of one piece, write its extension as
\begin{equation}\label{glob:patch-extension}
 F_0(y)=F_0(x,t)
 =\int_{\R^2}f(\xi)\chi_1^{\nu_0}(\xi)
       e^{2\pi i(x\cdot\xi+t\phi(\xi))}\,d\xi,
 \qquad \|f\|_\infty+\|f\|_2\le C.
\end{equation}
Here $f$ is supported in a fixed compact frequency patch. The smooth
surface-area density and division by $\chi_1^{\nu_0}$ have been absorbed
into $f$. All constants in this section may depend on these finitely many
fixed choices, but are independent of the original bounded measurable
data. In particular, no differentiability of $f$ is assumed.

\Needspace{4\baselineskip}
\begin{lemma}[Sparse cubic estimate]\label{glob:sparse-cubic}
There are constants $C_0>1$ and $\mu_0>0$ such that the following holds.
Let $0<\mu\le\mu_0$, let $R\ge\mu^{-1}$, and let $\mathcal B$ be a
collection of at most $\mu^{-8}$ balls of radius $R$ whose distinct
centers have distance at least $(R/\mu)^{C_0}$. Then, for every
$\epsilon>0$,
\begin{equation}\label{eq:source-42}
             \sum_{B\in\mathcal B}\int_B |F_0(y)|^3\,dy
                    \le C_\epsilon R^\epsilon.
\end{equation}
The constants are uniform over the data in
\eqref{glob:patch-extension} and all translations of the balls.
\end{lemma}

\begin{proof}
Choose a dyadic $N$ with $4R\le N^2<16R$, and put $\delta=N^{-1}$.
For a ball with center $b_B\in\R^3$, set
$a_B=b_B-(0,0,N^2/2)$. The coordinates
$y=a_B+N^2(z,t)$ place that ball inside
$\{|z|\le 1/4,\ 1/4\le t\le3/4\}$.
We shall also translate the physical sampling lattice by
\[
 u\in Q=[0,L_{\mathrm f}^{-1})^2\times[0,1).
\]
For each such $u$, define
\[
 f_{B,u}(\xi)=f(\xi)
       e^{2\pi i(a_B+u)\cdot(\xi,\phi(\xi))}.
\]
These functions have the same $L^\infty$ and $L^2$ bounds as $f$.
All estimates that follow are uniform in \(u\) and in the ball
locations.  Those locations enter only through this modulation of the
input; the retained packet coordinates stay in the same bounded
normalized region.

At factor $1$, retain the analysis coefficients of $f_{B,u}$ with
angular windows meeting the fixed frequency support and positions
\[
 z\in L_{\mathrm f}^{-1}\delta\Z^2,\qquad |z|\le Z_0,
\]
where $Z_0$ is a sufficiently large fixed constant. Denote the resulting
finite array by $d^{B,u}$. For each angle $\nu$, Fourier-series
Parseval gives
\begin{equation}\label{glob:angle-mass}
 \sum_z|d_{\nu,z}^{B,u}|^2
 \le (L_{\mathrm f}\delta)^2
      \int |f(\xi)\chi_\delta^\nu(\xi)|^2\,d\xi
 \le C\delta^4.
\end{equation}
The velocity centers $U_\nu=-\nabla\phi(\nu)$ are separated by at
least $c\delta$, since the Hessian of the extended phase is uniformly
definite. An ellipse $E$ whose radii satisfy $L\ge w\ge\delta$ thus
contains at most $C|E|\delta^{-2}$ such centers: its $O(\delta)$
enlargement has area comparable to $|E|$. Ignoring the positional
restriction in the definition of capacity, we obtain from
\eqref{eq:source-29} and \eqref{glob:angle-mass}
\begin{equation}\label{glob:initial-capacity}
 K(d^{B,u})
 \le C\delta^4\sup_E\frac{|E|}{W(E)}
 \le C\delta^4,
\end{equation}
because $|E|/W(E)$ is a constant multiple of $(w/L)^\kappa\le1$.

We next prove the estimate that saves the number of balls:
\begin{equation}\label{glob:joint-mass}
       \sum_{B\in\mathcal B}\sum_v|d_v^{B,u}|^2
                         \le C\delta^2.
\end{equation}
Let $T_B$ be the retained analysis operator acting on the unmodulated
function $f$, so that $T_Bf=d^{B,u}$, and let $T=(T_B)_{B\in\mathcal B}$.
The frame identity shows that each $T_B$ has operator norm at most
$L_{\mathrm f}\delta$. Hence the diagonal blocks of $TT^*$ have
operator norm at most $C\delta^2$.

An entry in an off-diagonal block is, up to complex conjugation, an
integral of the form
\begin{equation}\label{glob:gram-entry}
 \int A(\xi)e^{2\pi i(q\cdot\xi+s\phi(\xi))}\,d\xi,
 \qquad
 (q,s)=a_B-a_{B'}+N^2(z-z',0),
\end{equation}
where $A$ is a product of two $\delta$-scale windows supported in a fixed
compact set. The common translation $u$ cancels. The position offsets
have size $O(N^2)=O(R)$, so, for $B\ne B'$ and sufficiently small
$\mu_0$,
\[
 H:=|(q,s)|\ge c(R/\mu)^{C_0}.
\]
The amplitude and its derivatives have bounds polynomial in $R$.
There is a fixed exponent $a\ge0$, independent of $\kappa$, $\epsilon$,
and the ball locations, for which
\begin{equation}\label{glob:gram-decay}
 \left|\int A(\xi)e^{2\pi i(q\cdot\xi+s\phi(\xi))}\,d\xi\right|
                       \le C R^aH^{-1/2}.
\end{equation}
For completeness, write $\Psi(\xi)=q\cdot\xi+s\phi(\xi)$.
If $|s|$ is sufficiently small compared with $H$, then
$|\nabla\Psi|\ge cH$ on the fixed support and integration by parts
proves \eqref{glob:gram-decay}. Otherwise $|s|\ge cH$. Uniform
definiteness of the Hessian implies
\[
 |\nabla\Psi(\xi)-\nabla\Psi(\xi')|
                         \ge cH|\xi-\xi'|.
\]
Consequently the set $|\nabla\Psi|\le2H^{3/4}$ has area at most
$CH^{-1/2}$. Insert a smooth cutoff to this set. Its contribution is
bounded by $CH^{-1/2}\|A\|_\infty$. On the complementary region,
integrate by parts with the vector field
$\nabla\Psi/|\nabla\Psi|^2$. This field has size $O(H^{-3/4})$
and divergence $O(H^{-1/2})$; the derivative of the cutoff has size
$O(H^{1/4})$. One integration by parts therefore gives
$CR^aH^{-1/2}$, after increasing the fixed exponent $a$ to cover the
amplitude derivatives. This proves \eqref{glob:gram-decay}.

There are $O(N^2)$ retained angular labels and $O(N^2)$ retained
position labels per ball. Thus each off-diagonal row of $TT^*$ has at
most $C\mu^{-8}R^2$ entries. By \eqref{glob:gram-decay}, its absolute
row sum is at most
\[
 C\mu^{-8}R^{a+2}(R/\mu)^{-C_0/2}.
\]
Choose $C_0$ so large that $C_0/2\ge a+3$ and $C_0/2\ge8$.
This row sum is then $O(R^{-1})=O(\delta^2)$. The same holds for
columns, since the matrix is Hermitian. The Schur bound, together with
the diagonal block bound, proves $\|TT^*\|\le C\delta^2$ and hence
\eqref{glob:joint-mass}. This choice of $C_0$ is made once and does
not depend on the power loss $\epsilon$.

Apply the single-step packet map from factor $1$ to factor $N$ to
$d^{B,u}$. Retain only the terminal angle $\nu_0$ and positions in a
fixed bounded region containing $|z|\le1$, and denote the output by
$c^{B,u}$. At this terminal factor,
\[
       \theta=1,\qquad r=w=N^{-2},\qquad rw^2=N^{-6}.
\]
If the initial positions had not been truncated, the frame identity
would give exactly the samples
\[
       F_0\bigl(a_B+u+N^2(z,t_I)\bigr).
\]
The omitted positions contribute $O_M(N^{-M})$, uniformly on the
retained terminal region, for every fixed $M$.
Indeed, choose $Z_0$ larger than a fixed bound for the retained terminal
positions and the velocities on the frequency patch. If $|z'|>Z_0$,
then the phase in the synthesis--analysis matrix has, on a parent
$\delta$-cap, nonvanishing gradient of size at least
$cN^2(1+|z'|)$. Rescaling that cap and integrating by parts gives a
matrix bound $C_L(1+N|z'|)^{-L}$ for every $L$. There are $O(N^2)$
relevant parent angular labels, and their position mesh is comparable
to $N^{-1}$. Thus, for $L>1$, the squared matrix bounds sum to
\[
 \sum_{\nu\text{ relevant}}\sum_{z':\,|z'|>Z_0}(1+N|z'|)^{-2L}
                         \le C_LN^{4-2L}.
\]
Write $\widetilde d=\mathsf P_{1,[0,1)}f_{B,u}$ for the full initial
analysis array. The frame identity gives
$\sum_v|\widetilde d_v|^2\le C\delta^2$, so Cauchy--Schwarz bounds
each omitted contribution by $C_LN^{1-L}$. Taking $L\ge M+1$ proves
the claimed error, using only smoothness of the windows and the phase.

The supports just described belong to one fixed polynomial complexity
range. Theorem~\ref{pkt:main}, \eqref{eq:source-30}, and
\eqref{glob:initial-capacity} therefore give, for every fixed
$0<\kappa<1/10$ and $\eta>0$,
\begin{align}
 N^{-6}\sum_v|c_v^{B,u}|^3
 &\le \mathcal A(c^{B,u})^3
 \le C_{\kappa,\eta}N^{10\kappa+\eta}G(d^{B,u})^3 \notag\\
 &\le C_{\kappa,\eta}N^{10\kappa+\eta}
                 \delta^4\sum_v|d_v^{B,u}|^2.
 \label{glob:sparse-sampling}
\end{align}
In the last line we used $G(d)^3=\delta^2(\sum|d_v|^2)K(d)^{1/2}$
at the initial factor. Summing \eqref{glob:sparse-sampling} and using
\eqref{glob:joint-mass}, the physical sample sum is at most
\[
 C_{\kappa,\eta}N^6N^{10\kappa+\eta}\delta^4\delta^2
                    =C_{\kappa,\eta}N^{10\kappa+\eta}.
\]
The sampling errors are harmless: there are at most
$C\mu^{-8}N^6\le CN^{22}$ samples, and $M$ may be chosen arbitrarily
large. The inequality $|a+b|^3\le4(|a|^3+|b|^3)$ absorbs their total
cubed contribution.

The terminal physical lattice is
$\Lambda=L_{\mathrm f}^{-1}\Z^2\times\Z$, with fundamental cell $Q$.
The retained lattice points include each $\lambda\in\Lambda$ for which
$\lambda+Q$ meets $B-a_B$, since the ball lies in the interior
normalized region fixed above.
Consequently,
\[
 \sum_{B\in\mathcal B}\int_B|F_0(y)|^3\,dy
 \le\int_Q\sum_{B\in\mathcal B}\sum_{v\text{ retained}}
   |F_0(a_B+u+N^2(z_v,t_{I_v}))|^3\,du.
\]
Apply the uniform sampled bound to the right-hand side. Integration
over $Q$ tiles physical space by $\Lambda$, so it introduces no
$N$-dependent density factor. Finally choose
$\kappa$ and $\eta$ sufficiently small that
$(10\kappa+\eta)/2\le\epsilon$. Since $N^2\asymp R$, this proves
\eqref{eq:source-42}.
\end{proof}

\subsection{A finite list of peaks}

We now use the geometry of the original sphere to obtain a global
estimate at a larger exponent.  The convolution calculation keeps the
normalization of surface measure explicit and applies directly to
bounded measurable complex data.

\begin{lemma}[An elementary fourth-power bound]\label{glob:fourth-power}
For every bounded measurable $h:S^2\to\C$,
\begin{equation}\label{glob:l4-bound}
       \|Eh\|_{L^4(\R^3)}^4\le32\pi^3\|h\|_\infty^4.
\end{equation}
Moreover $Eh$ and its first derivatives are uniformly bounded by
constants times $\|h\|_\infty$.
\end{lemma}

\begin{proof}
The convolution of surface measure with itself has density
\begin{equation}\label{glob:sphere-convolution}
       (\sigma*\sigma)(y)=\frac{2\pi}{|y|}
                               \mathbf1_{\{0<|y|<2\}}.
\end{equation}
To verify the constant, use rotational invariance and fix
$\omega\in S^2$. Under surface area measure in the second variable,
$s=\omega\cdot\nu$ has density $2\pi$ on $[-1,1]$, and
$|\omega+\nu|=(2+2s)^{1/2}$. Thus, for a radial test function $q$,
\[
 \int q(|y|)\,d(\sigma*\sigma)(y)
     =8\pi^2\int_0^2q(r)r\,dr,
\]
which is exactly the integral against the radial density in
\eqref{glob:sphere-convolution}. In particular,
\[
       \|\sigma*\sigma\|_2^2
       =\int_0^2\frac{4\pi^2}{r^2}\,4\pi r^2\,dr
       =32\pi^3.
\]
Let $\lambda=h\sigma$, a finite complex measure. Its convolution
satisfies the total-variation domination
\[
       |\lambda*\lambda|\le\|h\|_\infty^2(\sigma*\sigma),
\]
so it has an $L^2$ density of squared norm at most
$32\pi^3\|h\|_\infty^4$. The function $(Eh)^2$ is the Fourier
transform, with the positive phase convention, of this convolution.
Plancherel proves \eqref{glob:l4-bound}. Finally, differentiating the
defining integral is justified by boundedness of the frequency support
and gives, for example,
\[
       \|Eh\|_\infty\le4\pi\|h\|_\infty,
       \qquad \|\nabla Eh\|_\infty\le8\pi^2\|h\|_\infty.
\]
\end{proof}

Each function $F_0$ in \eqref{glob:patch-extension}, in its rotated
coordinates, is the extension of a uniformly bounded density supported
on one of the original sphere pieces. Lemma~\ref{glob:fourth-power}
therefore supplies uniform $L^4$, $L^\infty$, and Lipschitz bounds for
$F_0$.

\begin{lemma}[Initial peak count]\label{glob:initial-peaks}
For $0<\mu\le1$, the number of unit lattice cubes $Q'$ satisfying
$\sup_{Q'}|F_0|\ge\mu$ is at most $C\mu^{-7}$.
One may select a point $x_{Q'}\in Q'$ in each such cube with
$|F_0(x_{Q'})|\ge\mu/2$.
\end{lemma}

\begin{proof}
Take the unit cubes to be half-open, and select such a point in every
cube of an arbitrary finite subcollection. A uniform Lipschitz bound
shows that $|F_0|\ge\mu/4$ on the ball of radius $c\mu$ about each
selected point, where $c>0$ is fixed and small. Each ball contributes
at least $c'\mu^7$ to the integral of $|F_0|^4$. These balls have
bounded overlap, since their radii are bounded and there is one center
in each of distinct unit lattice cubes. Lemma~\ref{glob:fourth-power}
bounds the size of the finite subcollection by $C\mu^{-7}$. Since
the same bound holds for every finite subcollection, the whole set of
such cubes is finite and has the asserted size.
\end{proof}

\subsection{Sparse covering and reproduction}

The covering lemma converts the initial polynomial peak count into
sparse families whose number is an arbitrarily small power of
\(\mu^{-1}\).  Their radii remain polynomial in \(\mu^{-1}\), with an
exponent fixed once \(\alpha\) is chosen.  This is the information
needed to apply the sparse cubic estimate and then choose its loss.

\begin{lemma}[Sparse covering]\label{glob:sparse-covering}
Fix $C_0>1$ and $A>0$. For every $0<\alpha<1$ there is a constant
$C_\alpha$ with the following property. If $\mu>0$ is sufficiently
small and $X$ is a finite indexed list of at most $A\mu^{-7}$ points
in $\R^3$, then $X$ is covered by at most
$C_\alpha\mu^{-\alpha/2}$ families of balls. Each family has a common
radius $\rho$ with
\begin{equation}\label{glob:covering-radii}
           \mu^{-1}\le\rho,\qquad 2\rho\le\mu^{-C_\alpha},
\end{equation}
and its distinct centers are separated by at least
$(2\rho/\mu)^{C_0}$. Each family has at most $\#X$ balls.
\end{lemma}

\begin{proof}
Put $\beta=\alpha/4$, choose an integer $J>8/\beta$, and set
$D=2C_0+2$. For sufficiently small $\mu$, we have $\#X\le\mu^{-8}$.
Define
\[
 R_0=\mu^{-1},\qquad R_{j+1}=(R_j/\mu)^D
              \quad(0\le j<J),
 \qquad n_j(x)=\#\bigl(X\cap B(x,R_j)\bigr),
\]
where balls are open and counts include the multiplicities of the
indexed list. For every point $x$ of the list, some $j<J$ satisfies
\begin{equation}\label{glob:slow-growth}
                n_{j+1}(x)\le\mu^{-\beta}n_j(x).
\end{equation}
Otherwise $n_J(x)>\mu^{-J\beta}n_0(x)\ge\mu^{-J\beta}$ would
contradict $\#X\le\mu^{-8}$.

Assign each point one such index $j$, and group the points further
according to $2^k\le n_j(x)<2^{k+1}$. There are
$O_\alpha(\log(1/\mu))$ groups. In each group choose a maximal
disjoint subcollection of the radius-$R_j$ balls centered at its
points. The triples of the selected balls cover that group.

Join two selected centers by an edge when their distance is less than
$R_{j+1}/2$. For a selected center $a$, the radius-$R_j$ balls
centered at its neighbors are disjoint and lie in $B(a,R_{j+1})$,
since $R_j\le R_{j+1}/2$. Each such ball captures at least $2^k$
indices of the whole list, whereas \eqref{glob:slow-growth} gives
$n_{j+1}(a)<2^{k+1}\mu^{-\beta}$. The graph therefore has degree
at most $2\mu^{-\beta}$. Greedy coloring splits the selected balls
into $O(\mu^{-\beta})$ colors. Including all the groups, the number
of families is at most
\[
 C_\alpha\mu^{-\beta}\log(1/\mu)
                         \le C_\alpha\mu^{-\alpha/2}.
\]

Use radius $\rho=3R_j$ for each covering family. Centers in one
color have distance at least $R_{j+1}/2$, and, for sufficiently small
$\mu$,
\[
       \frac{R_{j+1}}2
       =\frac12(R_j/\mu)^{2C_0+2}
       \ge(6R_j/\mu)^{C_0}
       =(2\rho/\mu)^{C_0}.
\]
Finally, write $R_j=\mu^{-a_j}$, so that
$a_0=1$ and $a_{j+1}=D(a_j+1)$. The finite number $J$ depends only
on $\alpha$ and $C_0$. Enlarging $C_\alpha$ to at least
$a_{J-1}+1$ gives $6R_j\le\mu^{-C_\alpha}$ for small $\mu$,
which proves \eqref{glob:covering-radii}. The selected balls are
centered at points of the list, so every family has at most $\#X$
balls.
\end{proof}

\Needspace{4\baselineskip}
\begin{proposition}[Distributional bound]\label{glob:distribution}
For every $0<\alpha<1$ and all sufficiently small $\mu>0$,
\begin{equation}\label{glob:distribution-bound}
       |\{y\in\R^3:|F_0(y)|>\mu\}|
                            \le C_\alpha\mu^{-3-\alpha}.
\end{equation}
\end{proposition}

\begin{proof}
Let $X=(x_{Q'})$ be the finite list supplied by
Lemma~\ref{glob:initial-peaks}. It suffices to prove
$\#X\le C_\alpha\mu^{-3-\alpha}$, since the level set in
\eqref{glob:distribution-bound} is contained in the corresponding
unit cubes.

The Fourier transform of $F_0$, as a tempered distribution, is
supported in a fixed compact set. Choose a Schwartz function
$\psi$ whose Fourier transform is $1$ on a neighborhood of that
set. Then $F_0=\psi*F_0$; the convolution is absolutely convergent
and the identity holds pointwise because $F_0$ is bounded and
continuous. Hölder's inequality yields
\begin{equation}\label{glob:reproduction}
 |F_0(x)|^3
 \le\|\psi\|_1^2
        \int_{\R^3}|\psi(x-y)|\,|F_0(y)|^3\,dy.
\end{equation}
Apply Lemma~\ref{glob:sparse-covering}, with the separation exponent
from Lemma~\ref{glob:sparse-cubic}, and assign every index of $X$
to one ball in one covering family. Suppose the assigned covering
ball has radius $\rho$. Truncating the integral in
\eqref{glob:reproduction} to $|x-y|\le\rho$ costs at most
\[
       C\|F_0\|_\infty^3
                \int_{|z|>\rho}|\psi(z)|\,dz
                    \le C\rho^{-4}\le C\mu^4.
\]
Since $|F_0(x)|\ge\mu/2$, this error is absorbed into the
left-hand side for sufficiently small $\mu$. The truncated region
lies in the doubled assigned covering ball.

There is a uniform bound
\begin{equation}\label{glob:kernel-overlap}
           \sup_{y\in\R^3}\sum_{x\in X}|\psi(x-y)|\le C_\psi.
\end{equation}
Indeed, at most one selected point lies in each unit lattice cube,
and Schwartz decay bounds the sum by the convergent lattice sum of
a sufficiently high inverse power of distance. In particular no
pairwise separation assumption on the selected points is needed.

Let $X_{\mathcal F}$ denote the indices assigned to one covering
family $\mathcal F$, of common radius $\rho$. Summing the truncated
form of \eqref{glob:reproduction} and using
\eqref{glob:kernel-overlap}, we obtain
\[
       \mu^3\#X_{\mathcal F}
       \le C\sum_{B\in\mathcal F}\int_{2B}|F_0(y)|^3\,dy
       \le C_\epsilon(2\rho)^\epsilon.
\]
For the last inequality, the doubled balls satisfy all the
hypotheses of Lemma~\ref{glob:sparse-cubic}: their radius is at
least $\mu^{-1}$, their separation is the required one, and their
number is at most $\#X\le\mu^{-8}$ after reducing the threshold
for $\mu$.

Fix $\alpha$ first, and then choose $\epsilon>0$ so small that
$C_\alpha\epsilon\le\alpha/2$, using the exponent in
\eqref{glob:covering-radii}. Thus
\[
       \mu^3\#X_{\mathcal F}\le C_\alpha\mu^{-\alpha/2}.
\]
There are at most $C_\alpha\mu^{-\alpha/2}$ families. Each index
was assigned exactly once, and summing proves
$\#X\le C_\alpha\mu^{-3-\alpha}$, as required.
\end{proof}

\begin{proof}[Proof of Theorem~\ref{thm:main}]
Fix $p>3$ and choose $0<\alpha<\min\{1,p-3\}$.
Let $\mu_*>0$ be smaller than the thresholds required above for this
choice of $\alpha$.
By the layer-cake identity and Proposition~\ref{glob:distribution},
the contribution of sufficiently small thresholds is bounded by
\[
 p\int_0^{\mu_*}\mu^{p-1}
      |\{|F_0|>\mu\}|\,d\mu
 \le C_{p,\alpha}\int_0^{\mu_*}\mu^{p-4-\alpha}\,d\mu
 <\infty.
\]
For larger thresholds, the uniform $L^4$ estimate gives
$|\{|F_0|>\mu\}|\le C\mu^{-4}$, and the level set is empty
above the uniform $L^\infty$ bound. Their contribution is therefore
finite as well. We conclude that $\|F_0\|_p\le C_p$.

There are only finitely many graph pieces in
\eqref{glob:patch-extension}. Rotations preserve Lebesgue measure,
so the $L^p$ triangle inequality gives $\|Eg\|_p\le C_p$ under
the normalization $\|g\|_\infty=1$. Multiplication by
$\|g\|_\infty$ proves the result for arbitrary nonzero bounded
measurable complex data. When the essential supremum is zero, the
extension is identically zero. All constructions and estimates
depend only on the almost-everywhere class of the data, which
completes the assertion of the theorem.
\end{proof}
